\begin{frame}{Contents}
\protect\hypertarget{in-queste-diapositive}{}

\begin{itemize}
\item
  Exercise 4
\item
  Exercise 1
\item
  Exercise 9
\item
  Exercise 10
\item
  Exercise 11
\end{itemize}

\end{frame}

\begin{frame}{Exercise 4}
What is the short-run effect on the exchange rate of an increase in the level of GNP in real terms, given exchange rate expectations?
\end{frame}

\begin{frame}{Exercise 1}
Suppose that there is a reduction in the aggregate real demand for money, i.e. a negative shift in the aggregate real demand function for money. Explain the short- and long-term effects on the exchange rate, interest rate and price level.
\end{frame}

\begin{frame}{Exercise 9}

In our discussion of short-run exchange rate overshooting, we assumed that real output was given. Instead, we assume that in the short run an increase in the money supply increases the level of real production (a hypothesis that will be justified in Chapter 6). How does this affect the over-reaction of the exchange rate due to the increase in the money supply? Is it possible for the exchange rate to show undershooting?
\end{frame}

\begin{frame}{Exercise 10}

Figure 3.2 shows that the short-term interest rate in Japan has gone through periods where it has been close to or even zero. Is the fact that the yen interest rates in the figure never turned negative just a coincidence, or can you find some explanation for why interest rates are constrained from the bottom to zero?
\end{frame}

\begin{frame}{Exercise 11}

How might a zero interest rate complicate the task of monetary policy?
\end{frame}

